\par\Needspace{6\baselineskip}\medskip
\hypertarget{native-204-article152}{}\pdfbookmark[2]{152.}{article-152}
\begin{center}152.\end{center}
Nachdem die vorbereitenden Rechnungen so abgeschlossen sind, gehen wir zur ersten Hypothese über. Das (unkorrigierte) Zeitintervall zwischen der zweiten und dritten Beobachtung beträgt $9{,}971\,192$ Tage, zwischen der ersten und zweiten $11{,}963\,241$ Tage. Die Logarithmen dieser Zahlen sind $0{,}998\,7471$ und $1{,}077\,8489$. Daraus folgt $\log\theta=9{,}234\,3285$, $\log\theta''=9{,}313\,4303$. Für die \textit{erste Hypothese} setzen wir daher
\[
\begin{aligned}
x&=\log P=0{,}079\,1018\\
y&=\log Q=8{,}547\,7588.
\end{aligned}
\]
Daraus ergibt sich $P=1{,}199\,7804$, $P+a=1{,}554\,1396$, $P+d=-0{,}162\,7248$.
\SourceLeaf{205}
\begin{center}
\begin{tabular}{@{}ll@{}}
\makebox[18em][l]{$\log e$\hfill .\,.\,.\,.\,.\,.\,.\,.\,.\,.} & $8{,}392\,9518_n$\\[.3em]
\makebox[18em][l]{$\log(P+a)$\hfill .\,.\,.\,.\,.} & $0{,}191\,4900$\\[.3em]
\makebox[18em][l]{C. $\log(P+d)$\hfill .\,.\,.} & $0{,}788\,5463_n$\\[.3em]
\hline
\makebox[18em][l]{$\log\operatorname{tang}\omega$\hfill .\,.\,.\,.\,.} & $9{,}372\,9881$\\[.3em]
\hline
\end{tabular}
\end{center}
\noindent daraus $\omega=+13^0\,16'\,51''{,}89$, $\omega+\sigma=+13^0\,40'\,5''{,}01$.
\begin{center}
\begin{tabular}{@{}ll@{}}
\makebox[18em][l]{$\log Q$\hfill .\,.\,.\,.\,.\,.\,.\,.\,.} & $8{,}547\,7588$\\[.3em]
\makebox[18em][l]{$\log c$\hfill .\,.\,.\,.\,.\,.\,.\,.\,.\,.} & $2{,}690\,7847$\\[.3em]
\makebox[18em][l]{$\log\sin\omega$\hfill .\,.\,.\,.\,.\,.\,.} & $9{,}361\,2147$\\[.3em]
\hline
\makebox[18em][l]{$\log Qc\sin\omega$\hfill .\,.\,.} & $0{,}599\,7582$\\[.3em]
\end{tabular}
\end{center}
\noindent Nach wenigen Versuchen findet man, dass die Gleichung $Qc\sin\omega\sin z^4=\sin(z-13^0\,40'\,5''{,}01)$ durch den Wert $z=14^0\,35'\,4''{,}90$ erfüllt wird. Daraus folgt $\log\sin z=9{,}401\,0744$, $\log r'=0{,}325\,1340$. Außer dieser Lösung lässt die Gleichung drei weitere zu, nämlich
\[
\begin{aligned}
z&=32^0\quad2'\,28''\\
z&=137\quad27\quad59\\
z&=193\quad4\quad18.
\end{aligned}
\]
Die dritte ist zu verwerfen, weil $\sin z$ negativ wird, die zweite, weil $z$ größer als $\delta'$ ist. Die erste entspricht der in Art. 142 besprochenen Näherung an die Erdbahn.

Weiter erhalten wir nach Art. 143
\begin{center}
\begin{tabular}{@{}ll@{}}
\makebox[18em][l]{$\log\dfrac{R'\sin\delta'}{b}$\hfill .\,.\,.\,.\,.\,.\,.\,.\,.} & $9{,}864\,8551$\\[.3em]
\makebox[18em][l]{$\log(P+a)$\hfill .\,.\,.\,.\,.\,.\,.\,.\,.} & $0{,}191\,4900$\\[.3em]
\makebox[18em][l]{C. $\log\sin(z-\sigma)$\hfill .\,.\,.\,.\,.} & $0{,}610\,3578$\\[.3em]
\hline
\makebox[18em][l]{$\log\dfrac{n'r'}{n}$\hfill .\,.\,.\,.\,.\,.\,.\,.\,.\,.\,.} & $0{,}666\,7029$\\[.3em]
\makebox[18em][l]{$\log P$\hfill .\,.\,.\,.\,.\,.\,.\,.\,.\,.\,.\,.\,.} & $0{,}079\,1018$\\[.3em]
\hline
\makebox[18em][l]{$\log\dfrac{n'r'}{n''}$\hfill .\,.\,.\,.\,.\,.\,.\,.\,.\,.\,.} & $0{,}587\,6011$\\[.8em]
\hline
\end{tabular}
\end{center}
\begin{center}
\begin{tabular}{@{}ll@{}}
\multicolumn{2}{@{}l@{}}{$z+A'D-\delta'=z+199^0\,47'\,1''{,}51=214^0\,22'\,6''{,}41$;}\\[.3em]
\makebox[18em][l]{$\log\sin$\hfill .\,.\,.\,.} & $9{,}751\,6736_n$\\[.3em]
\multicolumn{2}{@{}l@{}}{$z+A'D''-\delta'=z+188\quad54\quad32{,}94=203\quad29\quad37{,}84$;}\\[.3em]
\makebox[18em][l]{$\log\sin$\hfill .\,.\,.\,.} & $9{,}600\,5923_n$.\\[.3em]
\end{tabular}
\end{center}
\noindent Daraus ergibt sich $\log p=9{,}927\,0735_n$, $\log p''=0{,}022\,6459_n$ und dann $\log q=0{,}293\,0977_n$, $\log q''=0{,}258\,0086_n$, woraus folgt
\SourceLeaf{206}
\begin{center}
\begin{tabular}{@{}ll@{}}
$\zeta=203^0\,17'\,31''{,}22$ & $\log r=0{,}330\,0178$\\[.3em]
$\zeta''=210\quad10\quad58{,}88$ & $\log r''=0{,}321\,2819$.\\[.3em]
\end{tabular}
\end{center}
\noindent Schließlich erhalten wir nach Art. 144
\[
\begin{aligned}
\tfrac12(u''+u)&=205^0\,18'\,10''{,}53\\
\tfrac12(u''-u)&=-3\quad14\quad2{,}02\\
f'&=3\quad48\quad14{,}66
\end{aligned}
\]
\begin{center}
\begin{tabular}{@{}ll@{}}
\makebox[18em][l]{$\log\sin2f'$\hfill .\,.\,.\,.\,.\,.} & $9{,}121\,8791$\\[.3em]
\makebox[18em][l]{$\log r$\hfill .\,.\,.\,.\,.\,.\,.\,.\,.\,.} & $0{,}330\,0178$\\[.3em]
\makebox[18em][l]{C. $\log\dfrac{n'r'}{n}$\hfill .\,.\,.\,.\,.} & $9{,}333\,2971$\\[.8em]
\hline
\makebox[18em][l]{$\log\sin2f$\hfill .\,.\,.\,.\,.\,.} & $8{,}785\,1940$\\[.3em]
\multicolumn{2}{@{}l@{}}{$2f=3^0\,29'\,46''{,}03$}\\[.3em]
\end{tabular}
\end{center}
\begin{center}
\begin{tabular}{@{}ll@{}}
\makebox[18em][l]{$\log\sin2f'$\hfill .\,.\,.\,.\,.\,.} & $9{,}121\,8791$\\[.3em]
\makebox[18em][l]{$\log r''$\hfill .\,.\,.\,.\,.\,.\,.\,.\,.} & $0{,}321\,2819$\\[.3em]
\makebox[18em][l]{C. $\log\dfrac{n'r'}{n''}$\hfill .\,.\,.\,.\,.} & $9{,}412\,3989$\\[.8em]
\hline
\makebox[18em][l]{$\log\sin2f''$\hfill .\,.\,.\,.\,.\,.} & $8{,}855\,5599$\\[.3em]
\multicolumn{2}{@{}l@{}}{$2f''=4^0\,6'\,43''{,}28$}\\[.3em]
\end{tabular}
\end{center}
\noindent Die Summe $2f+2f''$ unterscheidet sich hier von $2f'$ nur um $0''{,}01$.

Um nun die Zeiten wegen der Aberration zu korrigieren, berechnet man die Abstände $\rho,\rho',\rho''$ nach den Formeln in Art. 145 und multipliziert anschließend die Zeit $493^{\mathrm{s}}$ oder $0^{\mathrm{d}}{,}005\,706$ mit diesen Abständen. Die Rechnung lautet
\begin{center}
\begin{tabular}{@{}ll@{}}
\makebox[18em][l]{$\log r$\hfill .\,.\,.\,.\,.\,.\,.\,.\,.\,.\,.} & $0{,}33002$\\[.3em]
\makebox[18em][l]{$\log\sin(AD'-\zeta)$\hfill .\,.} & $9{,}23606$\\[.3em]
\makebox[18em][l]{C. $\log\sin\delta$\hfill .\,.\,.\,.\,.\,.} & $0{,}50080$\\[.3em]
\hline
\makebox[18em][l]{$\log\rho$\hfill .\,.\,.\,.\,.\,.\,.\,.\,.\,.\,.} & $0{,}06688$\\[.3em]
\makebox[18em][l]{$\log$ Konst.\hfill .\,.\,.\,.\,.\,.\,.} & $7{,}75633$\\[.3em]
\hline
\makebox[18em][l]{$\log$ der Reduktion\hfill .\,.} & $7{,}82321$\\[.3em]
Reduktion $=$ & $0{,}006\,656$\\[.3em]
\end{tabular}
\end{center}
\begin{center}
\begin{tabular}{@{}ll@{}}
\makebox[18em][l]{$\log r'$\hfill .\,.\,.\,.\,.\,.\,.\,.} & $0{,}32513$\\[.3em]
\makebox[18em][l]{$\log\sin(\delta'-z)$\hfill .\,.} & $9{,}48384$\\[.3em]
\makebox[18em][l]{C. $\log\sin\delta'$\hfill .\,.\,.\,.} & $0{,}27189$\\[.3em]
\hline
\makebox[18em][l]{$\log\rho'$\hfill .\,.\,.\,.\,.\,.\,.\,.} & $0{,}08086$\\[.3em]
 & $7{,}75633$\\[.3em]
\hline
 & $7{,}83719$\\[.3em]
 & $0{,}006\,874$\\[.3em]
\end{tabular}
\end{center}
\begin{center}
\begin{tabular}{@{}ll@{}}
\makebox[18em][l]{$\log r''$\hfill .\,.\,.\,.\,.\,.\,.\,.\,.\,.} & $0{,}32128$\\[.3em]
\makebox[18em][l]{$\log\sin(A''D'-\zeta'')$\hfill .} & $9{,}61384$\\[.3em]
\makebox[18em][l]{C. $\log\sin\delta''$\hfill .\,.\,.\,.\,.\,.} & $0{,}16464$\\[.3em]
\hline
\makebox[18em][l]{$\log\rho''$\hfill .\,.\,.\,.\,.\,.\,.\,.\,.\,.\,.} & $0{,}09976$\\[.3em]
 & $7{,}75633$\\[.3em]
\hline
 & $7{,}85609$\\[.3em]
 & $0{,}007\,179$\\[.3em]
\end{tabular}
\end{center}
\begin{center}
\small\setlength{\tabcolsep}{8pt}\renewcommand{\arraystretch}{.85}
\begin{tabular}{@{}cccc@{}}
Beobachtung & korrigierte Zeiten & Intervalle & Logarithmen\\[.5em]
I. & Okt. $5{,}451\,988$ & & \\
 & & $11^{\mathrm{d}}{,}963\,023$ & $1{,}077\,8409$\\
II. & $17{,}415\,011$ & & \\
 & & $9{,}970\,887$ & $0{,}998\,7339$\\
III. & $27{,}385\,898$ & & \\
\end{tabular}
\end{center}
\noindent Die korrigierten Logarithmen der Größen $\theta,\theta''$ sind daher $9{,}234\,3153$ und $9{,}313\,4223$. Beginnt man nun die Bestimmung der Bahnelemente aus $f,r',r'',\theta$, so ergibt sich $\log\eta=0{,}000\,2285$, ebenso aus $f'',r,r',\theta''$ der Wert $\log\eta''=0{,}000\,3191$. Wir übergehen diese Rechnung hier, da sie im Ersten Buch, Abschnitt III, ausführlich erläutert wurde.

\SourceLeaf{207}Zuletzt erhalten wir nach Art. 146
\begin{center}
\begin{tabular}{@{}ll@{}}
\makebox[18em][l]{$\log\theta''$\hfill .\,.\,.\,.\,.\,.\,.\,.\,.\,.} & $9{,}313\,4223$\\[.3em]
\makebox[18em][l]{C. $\log\theta$\hfill .\,.\,.\,.\,.\,.\,.\,.} & $0{,}765\,6847$\\[.3em]
\makebox[18em][l]{$\log\eta$\hfill .\,.\,.\,.\,.\,.\,.\,.\,.\,.} & $0{,}000\,2285$\\[.3em]
\makebox[18em][l]{C. $\log\eta''$\hfill .\,.\,.\,.\,.\,.\,.} & $9{,}999\,6809$\\[.3em]
\hline
\makebox[18em][l]{$\log P'$\hfill .\,.\,.\,.\,.\,.\,.\,.\,.} & $0{,}079\,0164$\\[.3em]
\end{tabular}
\end{center}
\begin{center}
\begin{tabular}{@{}ll@{}}
\makebox[18em][l]{$2\log r'$\hfill .\,.\,.\,.\,.\,.\,.\,.\,.} & $0{,}650\,2680$\\[.3em]
\makebox[18em][l]{C. $\log rr''$\hfill .\,.\,.\,.\,.\,.\,.} & $9{,}348\,7003$\\[.3em]
\makebox[18em][l]{$\log\theta\theta''$\hfill .\,.\,.\,.\,.\,.\,.\,.\,.} & $8{,}547\,7376$\\[.3em]
\makebox[18em][l]{C. $\log\eta\eta''$\hfill .\,.\,.\,.\,.\,.} & $9{,}999\,4524$\\[.3em]
\makebox[18em][l]{C. $\log\cos f$\hfill .\,.\,.\,.\,.} & $0{,}000\,2022$\\[.3em]
\makebox[18em][l]{C. $\log\cos f'$\hfill .\,.\,.\,.\,.} & $0{,}000\,9579$\\[.3em]
\makebox[18em][l]{C. $\log\cos f''$\hfill .\,.\,.\,.\,.} & $0{,}000\,2797$\\[.3em]
\hline
\makebox[18em][l]{$\log Q'$\hfill .\,.\,.\,.\,.\,.\,.\,.\,.} & $8{,}547\,5981$\\[.3em]
\end{tabular}
\end{center}
\noindent Die erste Hypothese liefert also $X=-0{,}000\,0854$, \mbox{$Y=-0{,}000\,1607$}.
